\documentclass[conference]{IEEEtran}
\IEEEoverridecommandlockouts
\usepackage{cite}
\usepackage{amsmath,amssymb,amsfonts}
\usepackage{algorithmic}
\usepackage{graphicx}
\usepackage{textcomp}
\usepackage{xcolor}
\usepackage{booktabs}
\usepackage{multirow}
\usepackage[hidelinks]{hyperref}

\def\BibTeX{{\rm B\kern-.05em{\sc i\kern-.025em b}\kern-.08em
    T\kern-.1667em\lower.7ex\hbox{E}\kern-.125emX}}

\begin{document}

\title{Evaluating Low-Light Image Enhancement for Object Detection on ExDARK with YOLOv11n}

\author{\IEEEauthorblockN{Muhammad Iqbal Rasyid}
\IEEEauthorblockA{\textit{Dept. of Informatics} \\
\textit{Telkom University}\\
Bandung, Indonesia \\
otachiking@student.telkomuniversity.ac.id}
\and
\IEEEauthorblockN{Gamma Kosala}
\IEEEauthorblockA{\textit{Dept. of Informatics} \\
\textit{Telkom University}\\
Bandung, Indonesia \\
gammakosala@telkomuniversity.ac.id}
}


\maketitle

\begin{abstract}
This study evaluates the impact of integrating Low-Light Image Enhancement (LLIE) as a preprocessing stage on the detection performance of YOLOv11n using the Exclusively Dark (ExDARK) dataset. Three state-of-the-art LLIE methods (HVI-CIDNet, RetinexFormer, and LYT-Net) are systematically compared against a raw-image baseline using no-reference image quality metrics (NIQE, BRISQUE), structural edge preservation (EPI), detection accuracy (mAP@0.5, Precision, Recall), and computational cost (GFLOPs, latency). Experiments on 7{,}363 ExDARK images reveal a counter-intuitive finding: the baseline consistently outperforms all LLIE scenarios in mAP@0.5 (0.5997) and Recall (0.5569), despite the substantial perceptual quality improvements achieved by the enhancement methods. Eigen-CAM analysis on the C2PSA layer indicates that LLIE artifacts distort the variance distribution of learned features, redirecting the network's attention toward high-frequency noise rather than discriminative object regions. These results suggest that YOLOv11n possesses sufficient feature extraction capability for low-light conditions, and that adding LLIE preprocessing tends to be counterproductive to accuracy while introducing an average latency overhead of 101--356~ms per image.
\end{abstract}

\begin{IEEEkeywords}
Low-Light Image Enhancement, Object Detection, EigenCAM, YOLOv11, ExDARK, HVI-CIDNet
\end{IEEEkeywords}

% ========

\section{Introduction}
Object detection algorithms have advanced rapidly over the past decade. This progression began with the original \textit{You Only Look Once} (YOLO) architecture \cite{redmon2016yolo} and has since evolved through YOLOv8 and YOLO-NAS, which offer a compelling balance between speed and accuracy as comprehensively reviewed by Terven et al. \cite{terven2023yoloreview}, culminating in the latest iterations \cite{sapkota2025yoloreview}. Although these systems deliver strong performance under well-lit conditions, detection accuracy degrades substantially in low-light environments. The core challenges include reduced visibility, severely diminished contrast, and elevated sensor noise levels that obscure the structural details of objects \cite{loh2019exdark,hong2021lowlight}. Addressing these limitations carries direct societal relevance: reliable low-light detection underpins safety-critical applications including autonomous driving, urban surveillance, and smart city infrastructure, where consistent performance across all lighting conditions directly impacts human safety and public welfare \cite{liang2022edgeyolo}.

To address these challenges, Low-Light Image Enhancement (LLIE) methods have been proposed as a preprocessing stage to recover image visibility. A wide range of techniques have been developed, from classical Retinex-based approaches to more sophisticated deep learning methods such as RetinexFormer \cite{cai2023retinexformer} and the lightweight LYT-Net \cite{brateanu2025lytnet}. Among recent innovations, the Color and Intensity Decoupling Network (CIDNet) operating in the Horizontal/Vertical-Intensity (HVI) color space offers a principled separation of intensity and chrominance components to reduce restoration artifacts \cite{yan2025hvi}. However, performance evaluation of LLIE methods has predominantly relied on image quality metrics such as Peak Signal-to-Noise Ratio (PSNR) and Structural Similarity (SSIM), while their effectiveness on downstream vision tasks such as object detection is rarely validated directly, leaving a critical gap between perceptual quality gains and real-world task performance \cite{wu2024llieeffect,liu2021benchmark}.

When image enhancement methods have been integrated into detection pipelines, prior evaluations were largely confined to older-generation architectures. Mpouziotas et al. \cite{mpouziotas2022lowlight} demonstrated that enhancement methods can improve the performance of models such as YOLOv3 and YOLOv4. Similarly, Peng et al. \cite{peng2024yolov5llie} showed that integrating LLIE with YOLOv5 yielded significant detection accuracy gains in low-light settings. However, the feature extraction mechanisms in these earlier architectures are substantially more limited than those in current-generation detectors. In parallel, the evolution of compact detectors from YOLO Nano \cite{wong2019yolonano} to the latest YOLOv11 Nano variant (YOLOv11n) \cite{khanam2024yolov11} has substantially expanded their architectural capacity, yet whether these advances have rendered LLIE preprocessing unnecessary for such models, with its associated computational costs, remains largely uninvestigated.

To this end, this study systematically evaluates the effect of state-of-the-art LLIE preprocessing on the performance of YOLOv11n using the ExDARK dataset. Specifically, it makes four contributions. First, it assesses image enhancement quality using no-reference and structural metrics applicable to unpaired real-world data. Second, it presents a comparative analysis of detection performance across LLIE-enhanced scenarios. Third, it quantifies the computational overhead and latency penalty introduced by each enhancement pipeline. Fourth, it applies EigenCAM on the C2PSA pre-head feature layer to provide qualitative diagnostic insight into how LLIE artifacts distort the detector's spatial attention prior to classification.


% =======


\section{Related Work}
\subsection{Low-Light Image Enhancement and Image Quality Metrics}
The primary objective of LLIE algorithms is to recover visibility and color fidelity from images degraded by insufficient illumination \cite{liu2021benchmark}. State-of-the-art deep learning approaches such as RetinexFormer \cite{cai2023retinexformer} and the lightweight LYT-Net \cite{brateanu2025lytnet} have demonstrated strong capabilities in global illumination restoration. Meanwhile, CIDNet, which operates in the HVI color space, introduces a principled decoupling of intensity and chrominance components to prevent cross-channel noise amplification \cite{yan2025hvi}.

To assess the effectiveness of these methods, most LLIE studies rely on full-reference metrics such as PSNR and SSIM. However, these metrics cannot be applied objectively to real-world datasets like ExDARK that are inherently unpaired, lacking well-lit ground truth counterparts \cite{loh2019exdark,liu2021benchmark}. Accordingly, this study adopts no-reference metrics including the Natural Image Quality Evaluator (NIQE) \cite{mittal2013niqe} and the Blind/Referenceless Image Spatial Quality Evaluator (BRISQUE) \cite{mittal2012brisque}. Additionally, structural preservation of object boundaries is evaluated using the Edge Preservation Index (EPI) to ensure that illumination improvements do not compromise the spatial features critical for detection.

\subsection{Evolution of YOLO Architectures for Object Detection}
The YOLO family of algorithms has dominated real-time object detection owing to its architecturally efficient design since its first iteration \cite{redmon2016yolo,terven2023yoloreview}. Earlier generations such as YOLOv3 and YOLOv5 had limited feature extraction capacity, particularly for images with extremely low signal-to-noise ratios.

By contrast, the modern YOLOv11 architecture offers substantial improvements in both parameter efficiency and accuracy \cite{khanam2024yolov11}. For its ultra-compact Nano variant, designed specifically for edge devices with constrained computational budgets, floating-point operation (FLOPs) efficiency is particularly critical \cite{akavaram2025flops,dobrzycki2025freezing}. The integration of modules such as C3k2 and the C2PSA spatial attention mechanism enables even this lightweight model to extract raw gradient features effectively, despite operating within the strict latency budgets required for industrial deployment \cite{liang2022edgeyolo}.

\subsection{Impact of LLIE on Object Detection}
The majority of LLIE research focuses on improving visual quality for human perception without further validation on high-level machine vision tasks. For older-generation detectors with limited feature extraction capacity, applying LLIE as a preprocessing step has been shown empirically to improve detection accuracy \cite{wang2023yolov5lowlight,shovo2024lowlight}. As an alternative approach, some studies have modified the detector architecture itself, such as embedding spatial edge enhancement modules directly into the YOLOv8 backbone \cite{gong2025multiscale}.

Nevertheless, when LLIE is applied purely as external preprocessing, a fundamental disconnect exists between optimizing images for human perception and optimizing features for computational analysis. This was demonstrated through a comprehensive empirical study by Wu et al. \cite{wu2024llieeffect}, which evaluated LLIE effectiveness on modern detectors including RetinaNet, EfficientDet, and YOLOv7. Their results revealed that illumination recovery processes are prone to injecting high-frequency noise, artifacts, and semantic shifts that alter the original gradient distribution of the image. For highly capable detectors such as YOLOv11, LLIE artifacts may disrupt the network's attention focus. Rather than improving mAP, such preprocessing risks degrading overall detection accuracy while simultaneously burdening edge systems with disproportionate computational latency.

\section{Methodology}

\subsection{System Overview}
This study evaluates the end-to-end impact of LLIE preprocessing on object detection performance across three evaluation dimensions: image enhancement quality, detection accuracy, and total computational cost. The experimental pipeline is divided into four parallel scenarios. The first scenario serves as the baseline, feeding raw images directly into the detector. The second through fourth scenarios sequentially apply HVI-CIDNet, RetinexFormer, and LYT-Net as image enhancement modules prior to detection. An overview of the system design is shown in Fig.~\ref{fig:flowchart}.

\begin{figure}[htbp]
\centerline{\includegraphics[width=\linewidth]{flowchart_methodology.png}}
\caption{System overview showing the comparative pipeline between the raw-image baseline and LLIE-enhanced scenarios.}
\label{fig:flowchart}
\end{figure}

\subsection{Dataset and Preprocessing}
This study uses the ExDARK dataset, which was specifically designed for object detection under low-light conditions and has been widely adopted as a standard benchmark \cite{loh2019exdark}. The dataset spans 12 object categories across 10 illumination levels. Data splits follow the protocol established in the official ExDARK repository \cite{loh2019exdark}, yielding 3{,}000 images for training, 1{,}800 for validation, and 2{,}563 for testing.

Bounding box annotations are converted to the relative YOLO format, which encodes normalized center coordinates along with width and height ratios relative to the image dimensions \cite{padilla2021metrics,lin2014coco}. To meet the input requirements of YOLOv11 without introducing geometric distortion, all images are resized to 640$\times$640 pixels using letterbox resizing, which preserves the original aspect ratio by applying padding to the shorter dimension.

\subsection{Image Enhancement Inference}
Three LLIE methods are applied via zero-shot inference using weights pretrained on the LOL-v1 dataset: HVI-CIDNet \cite{yan2025hvi,hvicidnet2025}, RetinexFormer \cite{cai2023retinexformer,retinexformer2023}, and LYT-Net \cite{brateanu2025lytnet,lytnets2024}. Since ExDARK contains no paired well-lit references, this zero-shot protocol directly assesses each method's domain generalization capability without fine-tuning. All 7{,}363 images are enhanced, producing three fully enhanced dataset variants for experimental comparison.

\subsection{Detection Model Training}
The detector used in this study is YOLOv11n \cite{khanam2024yolov11}, the lowest-complexity variant of the YOLOv11 family. Selecting this variant serves to isolate the experimental variable: smaller-capacity models are more sensitive to input feature quality \cite{dobrzycki2025freezing}, so accuracy fluctuations more directly reflect the influence of LLIE rather than architectural redundancy. Its lightweight nature also partially offsets the computational overhead introduced by LLIE preprocessing.

Training is conducted identically across all four scenarios for up to 100 epochs with a batch size of 16. Early stopping is applied with a patience of 25 epochs, and all four models converged before reaching the maximum training budget. The AdamW optimizer \cite{loshchilov2017sgdr} is configured with an initial learning rate of 0.005, a final learning rate ratio of 0.01, and weight decay of 0.0005. A 5-epoch learning rate warmup stabilizes early gradient updates. Data augmentation includes Mosaic (1.0), Mixup (0.2), Copy-Paste (0.1), RandAugment, and Random Erasing (0.2), with Mosaic and Mixup disabled in the final 10 epochs to facilitate convergence on unmodified images.

\subsection{Evaluation Metrics and Interpretability}
Image enhancement quality is assessed using no-reference perceptual metrics NIQE \cite{mittal2013niqe} and BRISQUE \cite{mittal2012brisque}, alongside two spatial distribution metrics. Shadow Area (SA) quantifies the proportion of dark pixels below an intensity threshold of 30 (grayscale 0--255):
\begin{equation}
  \text{SA (\%)} = \frac{1}{N} \sum_{i=1}^{N} \mathbf{1}[I_i < 30] \times 100
  \label{eq:shadow}
\end{equation}
Root Mean Square (RMS) Contrast is defined as the standard deviation of pixel intensities:
\begin{equation}
  \text{RMS} = \sqrt{\frac{1}{N}\sum_{i=1}^{N}(I_i - \bar{I})^2}
  \label{eq:rms}
\end{equation}
where $I_i$ is the grayscale intensity of pixel $i$ and $\bar{I}$ is the mean intensity. Structural edge preservation is quantified by the Edge Preservation Index (EPI), computed as the cosine similarity between Sobel gradient magnitude vectors of the original ($G_o$) and enhanced ($G_e$) images:
\begin{equation}
  \text{EPI} = \frac{\sum_{x,y} G_e(x,y) \cdot G_o(x,y)}
                    {\sqrt{\sum_{x,y} G_e^2(x,y) \cdot \sum_{x,y} G_o^2(x,y)}}
  \label{eq:epi}
\end{equation}
A value of 1.0 indicates perfect structural preservation; lower values reflect edge degradation.

Detection performance follows the COCO evaluation standard \cite{padilla2021metrics,lin2014coco}, reporting mAP@0.5, mAP@0.5:0.95, per-class Average Precision (AP), Precision, and Recall. Computational cost is quantified via GFLOPs \cite{akavaram2025flops} and preprocessing latency (ms).

As a supplementary analysis, EigenCAM \cite{muhammad2020eigencam} is applied to the C2PSA layer of YOLOv11n (the final spatial attention layer before the Detection Head) to examine shifts in the detector's spatial attention across scenarios. EigenCAM extracts the first principal component (PC1) of the feature activation tensor via Singular Value Decomposition (SVD) in a class-agnostic and gradient-free manner, highlighting the regions that most strongly drive network attention.

\section{Results and Discussion}
\subsection{Image Enhancement}

This section analyzes the impact of LLIE preprocessing on both the perceptual quality and spatial structure of the images. Notably, all LLIE models operate in a zero-shot setting without additional pretraining on the target dataset. Fig.~\ref{fig:llie_samples} presents a visual comparison across six representative samples. In general, all methods successfully reveal previously obscured details through dynamic range expansion, as visible in samples (a), (c), and (d). However, the enhancement process introduces varying degrees of artifacts depending on the method.

HVI-CIDNet achieves the most aggressive brightness restoration but is susceptible to chromatic noise and color distortion, as observed in samples (b) and (d). In scenes with mixed lighting such as sample (e), this method exhibits highlight saturation leading to texture loss in bright regions, indicating a weakness in local tone mapping. RetinexFormer produces distinctive hue shift artifacts, an unnatural orange coloration in high-highlight regions (annotated by red arrow in samples (e) and (f)), likely caused by Retinex decomposition failure at high-gradient transitions. LYT-Net consistently exhibits the lowest noise levels visually and is free from significant color artifacts, while still revealing new objects and details in extremely dark images such as samples (b) and (c).

\begin{figure}[htbp]
    \centering
    \includegraphics[width=\columnwidth]{LLIE_6Samples.png}
    \caption{Visual comparison of LLIE results on six low-light samples: (a) vehicle exterior, (b) dim interior, (c) mixed-light with desk lamp, (d) dark room with minimal ambient light, (e) nighttime building exterior, (f) candlelit group scene. Columns: Raw, HVI-CIDNet, RetinexFormer, LYT-Net.}
    \label{fig:llie_samples}
\end{figure}
Table~\ref{tab:image_quality} presents the quantitative evaluation results. HVI-CIDNet achieves the best global illumination restoration, with the largest SA reduction (Eq.~\ref{eq:shadow}: 66.10 to 11.64) and the highest RMS Contrast (Eq.~\ref{eq:rms}: 61.02), supported by the best NIQE score indicating the most perceptually natural output. However, this optimization comes at the cost of low-level feature preservation: spatial noise ($\sigma$) increases across all methods, and the EPI (Eq.~\ref{eq:epi}) drops substantially from the baseline value of 1.00. LYT-Net achieves the best EPI among the three methods (0.83), consistent with its low-noise characteristics, while HVI-CIDNet records the lowest EPI (0.75), indicating the most aggressive structural edge disruption.

\begin{table}[htbp]
\centering
\caption{Image Quality Metrics Across LLIE Scenarios}
\label{tab:image_quality}
\resizebox{\columnwidth}{!}{%
\begin{tabular}{lcccc}
\toprule
\textbf{Metric} & \textbf{Baseline} & \textbf{HVI-CIDNet}
    & \textbf{RetinexFormer} & \textbf{LYT-Net} \\
\midrule
Shadow Area $\downarrow$   & 66.10 & \textbf{11.64} & 21.61 & \underline{13.65} \\
RMS Contrast $\uparrow$    & 35.17 & \textbf{61.02} & 56.49 & \underline{57.83} \\
NIQE $\downarrow$          & 5.18  & \textbf{3.92}  & \underline{3.94}  & 4.33  \\
BRISQUE $\downarrow$       & 31.14 & \underline{25.74} & \textbf{17.17} & 26.41 \\
Noise $\sigma$ $\downarrow$ & \textbf{4.52} & 7.43 & \underline{6.83} & 7.14 \\
EPI $\uparrow$             & \textbf{1.00} & 0.75 & 0.80 & \underline{0.83} \\
\bottomrule
\multicolumn{5}{l}{\footnotesize $\downarrow$ lower is better; $\uparrow$ higher is better. \textbf{Bold} = 1st; \underline{underline} = 2nd. SA in \%.}\\
\end{tabular}%
}
\end{table}

% ==============

\subsection{Object Detection}

The results in Table~\ref{tab:detection_metrics} reveal an empirical anomaly consistent with the findings of Wu et al.~\cite{wu2024llieeffect}: the baseline scenario maintains superiority in mAP@0.5 (0.5997) and Recall (0.5569), although LYT-Net Precision (0.6740) slightly exceeds the baseline (0.6644). This finding contrasts with prior work~\cite{wang2023yolov5lowlight,shovo2024lowlight} suggesting that LLIE preprocessing tends to improve detection accuracy. The discrepancy is likely attributable to the evolution of detector capabilities: YOLOv3 and YOLOv5, with their limited feature extraction capacity, indeed benefited from LLIE to clarify object pixels. In contrast, YOLOv11 incorporates C3k2, SPPF, and C2PSA spatial attention modules~\cite{khanam2024yolov11} that can extract raw gradient features without LLIE assistance, meaning that injected artifacts may instead disrupt the network's attention focus.

At the class level, LLIE impacts are non-uniform and class-dependent. Complex-boundary classes such as Car (0.6907 under LYT-Net vs. Baseline 0.7161) and Cat (0.3964 under CIDNet vs. Baseline 0.4418) consistently degrade across all LLIE methods, indicating sensitivity to edge distortion. However, several classes show clear improvement under LLIE: Bus (CIDNet: 0.8299 vs. Baseline: 0.7965), Chair (CIDNet: 0.5254 vs. Baseline: 0.5057), and Dog (LYT-Net: 0.5677 vs. Baseline: 0.5511). These classes share broad spatial extent and relatively simple boundaries, making them less sensitive to gradient noise and more responsive to improved visibility from illumination restoration. This class-level heterogeneity indicates that LLIE utility is object-class dependent rather than uniformly negative.

Qualitative evidence is presented in Fig.~\ref{fig:detection_6samples}. In sample (a), all scenarios detect Car and Bicycle, though LLIE scenarios tend to produce over-expanded bounding boxes due to dynamic range widening. In sample (b), both the Baseline and RetinexFormer fail to detect the cat in extreme darkness, while HVI-CIDNet produces a false positive Bottle from orange restoration artifacts; LYT-Net is the only scenario to detect the cat correctly. Sample (c) shows that LLIE improves visibility of a human figure that is difficult to distinguish in the raw image, yet the Baseline already detects the person despite the darkness, confirming that YOLOv11n can extract sufficient gradient features without explicit enhancement. In sample (d), dual occlusion reveals that CIDNet detects two chairs and one table while other scenarios recover only one or none, suggesting that LLIE-amplified noise can exacerbate detection errors under occlusion. In sample (e), despite HVI-CIDNet's aggressive brightening of the nighttime building exterior, all scenarios successfully detect the target object, illustrating that over-enhancement does not always impair detection. Sample (f) demonstrates that neither LLIE scenarios nor the Baseline are immune to false positives: HVI-CIDNet produces a double-detection of table bounding boxes, while the Baseline misidentifies a background region as a person; RetinexFormer delivers the qualitatively best performance on this sample.

\begin{table}[htbp]
\caption{Detection Performance Evaluation}
\label{tab:detection_metrics}
\centering
\resizebox{\columnwidth}{!}{%
\begin{tabular}{lcccc}
\toprule
\textbf{Metric} & \textbf{Baseline} & \textbf{CIDNet} & \textbf{RetinexFormer} & \textbf{LYT-Net} \\
\midrule
\textbf{mAP@0.5} $\uparrow$
  & \textbf{0.5997}
  & \underline{0.5914}
  & 0.5911
  & 0.5906 \\
\textbf{mAP@0.5:0.95} $\uparrow$
  & \textbf{0.3633}
  & \underline{0.3615}
  & 0.3563
  & 0.3554 \\
\textbf{Precision} $\uparrow$
  & \underline{0.6644}
  & 0.6657
  & 0.6600
  & \textbf{0.6740} \\
\textbf{Recall} $\uparrow$
  & \textbf{0.5569}
  & \underline{0.5497}
  & 0.5398
  & 0.5393 \\
\midrule
AP: Bicycle
  & \textbf{0.7443}
  & \underline{0.7374}
  & 0.7445
  & 0.7254 \\
AP: Boat
  & \textbf{0.6531}
  & \underline{0.6444}
  & 0.6308
  & 0.6298 \\
AP: Bottle
  & \underline{0.6417}
  & 0.6252
  & \textbf{0.6523}
  & 0.6435 \\
AP: Bus
  & \underline{0.7965}
  & \textbf{0.8299}
  & 0.8248
  & 0.8129 \\
AP: Car
  & \textbf{0.7161}
  & \underline{0.6980}
  & 0.6916
  & 0.6907 \\
AP: Cat
  & \textbf{0.4418}
  & 0.3964
  & \underline{0.4085}
  & 0.4034 \\
AP: Chair
  & \underline{0.5057}
  & \textbf{0.5254}
  & 0.5148
  & 0.4988 \\
AP: Cup
  & \underline{0.5579}
  & 0.5552
  & 0.5083
  & \textbf{0.5615} \\
AP: Dog
  & \underline{0.5511}
  & 0.5324
  & 0.5578
  & \textbf{0.5677} \\
AP: Motorbike
  & \underline{0.4687}
  & 0.4709
  & \textbf{0.4808}
  & 0.4715 \\
AP: People
  & \textbf{0.6831}
  & 0.6762
  & \underline{0.6780}
  & 0.6757 \\
AP: Table
  & \textbf{0.4358}
  & 0.3955
  & \underline{0.4071}
  & 0.4097 \\
\bottomrule
\end{tabular}%
}
{\footnotesize $\downarrow$ lower is better; $\uparrow$ higher is better. \textbf{Bold} = 1st; \underline{underline} = 2nd.}
\end{table}

\begin{figure}[htbp]
    \centering
    \includegraphics[width=\columnwidth]{Detection_6Samples_100.png}
    \caption{Detection predictions on the same six samples as Fig.~\ref{fig:llie_samples}. Columns: Ground Truth (GT), Baseline, HVI-CIDNet, RetinexFormer, LYT-Net.}
    \label{fig:detection_6samples}
\end{figure}
% ==============
\subsection{Computational Cost}

Table~\ref{tab:latency} presents the computational overhead introduced by each LLIE preprocessing pipeline, measured via enhancer latency $T_{\text{enhance}}$ and total pipeline GFLOPs. The baseline without LLIE serves as the reference point with zero enhancement overhead and a total of only 3.226 GFLOPs from the YOLO detector.

RetinexFormer incurs the highest latency and computational load (356.29~ms; 109.586 total GFLOPs), making it the least viable candidate for edge deployment. HVI-CIDNet and LYT-Net exhibit comparable enhancement latencies ($\approx$100~ms), but an important distinction exists: HVI-CIDNet processes 50.83 GFLOPs within that duration, indicating higher hardware utilization, while LYT-Net requires only 10.63 GFLOPs for similar latency, reflecting superior architectural efficiency and greater scalability on resource-constrained edge devices.

Regardless of these differences, all LLIE pipelines inherently exceed the standard real-time threshold of 33~ms per frame (30~FPS) typically required for edge-based detection systems \cite{liang2022edgeyolo}.

\begin{table}[htbp]
\centering
\caption{Pipeline Latency and Computational Cost}
\label{tab:latency}
\resizebox{\columnwidth}{!}{%
\begin{tabular}{lcc}
\toprule
\textbf{Scenario} & $T_{\text{enhance}}$ \textbf{(ms)}
    & \textbf{Total GFLOPs} \\
\midrule
Baseline      & 0.00   & 3.23  \\
HVI-CIDNet    & 100.53 & 54.06 \\
RetinexFormer & 356.29 & 109.59 \\
LYT-Net       & 101.67 & 13.86 \\
\bottomrule
\multicolumn{3}{l}{\footnotesize Total GFLOPs = Enhancer + YOLO GFLOPs (3.23), constant across scenarios.}\\
\end{tabular}%
}
\end{table}

% ==========

\subsection{Visual Interpretability}
EigenCAM \cite{muhammad2020eigencam} is applied to the C2PSA layer, the final spatial attention layer before the Detection Head in YOLOv11n, extracting the dominant variance direction (PC1) via Singular Value Decomposition (SVD). This class-agnostic and gradient-free method highlights spatial regions most strongly driving network attention, providing diagnostic insight into how LLIE artifacts alter the detector's internal representations.

EigenCAM operates at the feature level, not the prediction level. A functional separation exists between the backbone/neck that produces feature maps and the Detection Head that converts them into bounding boxes and scores \cite{zhou2016cam}. Consequently, saliency in the EigenCAM map does not directly correspond to detection success, which enables particularly rich diagnostic contrast across scenarios.

Two representative scenarios are illustrated in Fig.~\ref{fig:eigen_cam}. In case \textbf{(A)}, a cat in extreme darkness: S1 (Baseline) exhibits no discernible activation over the cat's body, reflecting insufficient salient gradient energy in the raw dark image. S2 (HVI-CIDNet) amplifies illumination but activates nearly the entire scene, with a dominant saliency hotspot toward an orange background object (a restoration artifact), causing the Detection Head to generate a false positive Bottle prediction. S3 (RetinexFormer) activates the cat's body region but also responds to background areas, resulting in a missed detection. S4 (LYT-Net) produces the most spatially focused activation over the cat's body, yielding the only correct detection in this scene.

In case \textbf{(B)}, a scene containing several Boats near the Eiffel Tower: all scenarios exhibit EigenCAM activation dominated by the tower, the most spatially prominent structure with large area and complex texture. As the Eiffel Tower is not a target class in ExDARK, this dominance illustrates PC1 capturing the globally maximal variance rather than class-relevant features. The Boats contribute minimal variance to PC1 and are not prominently highlighted. Nevertheless, only S4 (LYT-Net) successfully detects a Boat, while S2 produces a false detection, confirming that the Detection Head can respond to features not dominant in PC1, since lower principal components encoding small-object representations remain accessible to the full feature map \cite{zhou2016cam}.

\begin{figure}[htbp]
\centerline{\includegraphics[width=\columnwidth]{Eigen_3Sample_100_Fix.png}}
\caption{EigenCAM (PC1) activation maps on the C2PSA layer across four detection scenarios. Columns: (1) Raw, (2) HVI-CIDNet, (3) RetinexFormer, (4) LYT-Net. (A) Cat scene: S1 shows no activation on cat; S2 activates entire scene due to restoration artifacts. (B) Eiffel Tower scene: PC1 dominated by tower structure; only S4 detects a Boat.}
\label{fig:eigen_cam}
\end{figure}

\section{Conclusion}
This study has several limitations that should be acknowledged. The findings are specific to YOLOv11n and may not generalize to larger variants with greater feature extraction capacity. The LLIE models were applied zero-shot using weights pretrained on LOL-v1, and results may differ under domain-specific fine-tuning on ExDARK. The training pipeline employed minimal color-space augmentation (small HSV jitter without hard color jitter), which may limit the model's robustness to LLIE-induced color shifts; more aggressive chrominance augmentation could alter the sensitivity of these findings. Furthermore, EigenCAM visualizes only the first principal component (PC1) of the C2PSA feature tensor, leaving contributions from lower components unexamined.

Overall, the results demonstrate empirically that integrating LLIE algorithms (HVI-CIDNet, RetinexFormer, and LYT-Net) as preprocessing stages tends to be counterproductive to the detection performance of YOLOv11n under the tested conditions. Although these methods effectively reduce shadow areas (Eq.~\ref{eq:shadow}) and improve perceptual quality metrics, the restoration process degrades edge feature preservation, as reflected by EPI scores (Eq.~\ref{eq:epi}) declining from the baseline value of 1.00 to as low as 0.75 under HVI-CIDNet. EigenCAM analysis on the C2PSA layer reveals that LLIE artifacts redirect the network's attention toward background noise, potentially contributing to elevated false positive rates. At the same time, class-level analysis shows that LLIE provides measurable benefits for specific categories such as Bus, Chair, and Dog, indicating that its utility is class-dependent rather than uniformly detrimental. Furthermore, each LLIE module introduces a latency overhead of 101--356~ms per image, substantially exceeding real-time thresholds for edge deployment. The Baseline consistently outperforms all LLIE scenarios in mAP@0.5 and Recall, indicating that YOLOv11n possesses sufficient feature extraction capacity for low-light imagery without additional preprocessing. From a sustainable AI perspective, these findings advocate for leaner, single-stage pipelines that eliminate unnecessary computational overhead, supporting more energy-efficient and practically deployable vision systems in resource-constrained environments.

\bibliographystyle{IEEEtran}
\bibliography{referensi}


\end{document}